\documentclass[10pt,a4paper]{amsart}
\usepackage[margin=19mm]{geometry}
\usepackage{amsmath,amssymb,mathtools}
\usepackage[scr=rsfs,cal=euler]{mathalfa}
\usepackage{xfrac}
\usepackage[unicode,hidelinks,bookmarks=false]{hyperref}
\newcommand{\ie}{i.e.\ }
\newcommand{\cf}{cf.\ }
\newcommand{\resp}{respectively\ }
\newcommand{\Subsection}{\subsection}
\providecommand{\nobreakdash}{\nobreak\hbox{-}}
\setlength{\emergencystretch}{2em}
\multlinegap0pt
\usepackage{bm}
\usepackage{amssymb}
\newtheorem{proposition}{Proposition}
\newtheorem{theorem}{Theorem}
\renewcommand{\epsilon}{\varepsilon}
\renewcommand{\geq}{\geqslant}
\title{Height functions on singular surfaces parameterized by smooth maps $\mathcal{A}$-equivalent to $H_k$}
\author{Masaru Hasegawa}
\date{Source: arXiv:2608.22712v4 (CC BY 4.0)}
\begin{document}
\maketitle

\noindent\emph{Source and licence.} Height functions on singular surfaces parameterized by smooth maps $\mathcal{A}$-equivalent to $H_k$. Version arXiv:2608.22712v4 (CC BY 4.0), distributed by arXiv under the Creative Commons Attribution 4.0 International licence, http://creativecommons.org/licenses/by/4.0/, as recorded in the arXiv licence metadata for this exact version and shown on the abstract page https://arxiv.org/abs/2608.22712v4. Full licence notice: \texttt{licenses/arxiv-2608.22712-CC-BY-4.0.txt}.\par\smallskip

\noindent\textit{Editorial scope.} Counted unit: the article's theorem thm:branch (source0.tex lines 766--772). The article's complete original proof of it is delivered with it (source0.tex lines 773--806, one \texttt{proof} environment of 34 lines). No mathematics is written, repaired or re-derived: the statement and the proof are the article's own lines, in the article's own notation, and no retained line is substituted or re-flowed.  Retained uncounted as the source-local context the proof needs: lines 281--284; lines 589--613.  Nothing was omitted from any retained span, so the package carries no omission record.  No retained line contains a citation, so the wrapper carries no bibliography and no bibliography entry.  No retained line contains a figure environment, an image-inclusion command or drawing code, so no image is delivered and the package declares no asset dependency.  Editorial formatting only, in addition: an AMS \texttt{amsart} preamble that restates the article's own declarations and macros used by the retained lines, namely the environments \texttt{proposition}, \texttt{theorem} on separate counters, together with the standard packages those lines need.  The retained environments print in this document as Proposition 1, Theorem 1 in order of appearance, whereas the article prints the same items with the numbering of its own full document.  The complete native source of the article as distributed in the arXiv e-print for this exact version is bundled unchanged as \texttt{sources/208359/source0.tex}.
\begin{equation}
\label{eq:normal_form_2}
\left(u,\, u v + \sum_{i=3}^k \frac{b_i}{i!} v^i + O(u,v)^{k+1},\, \frac12 a_{2,0} u^2 + \sum_{m=3}^k\sum_{i+j=m} \frac{a_{i,j}}{i!j!} u^i v^j + O(u,v)^{k+1}\right), 
\end{equation}

\begin{proposition}
\label{prop:height}
Let $S$ be parameterized by $f$ in the form \eqref{eq:normal_form_2}.
\begin{enumerate}
\item 
$h_{\bm{w}_0}$ has an $A_1$-singularity at $(0,0)$ if and only if $\bm{w}_0$ is in the normal plane $($i.e., $x_0 = 0)$ but $\bm{w_0}$ is not the principal normal vector $($i.e., $\bm{w}_0 \ne \pm(0,0,1))$.
When this is the case, $H$ is an $\mathcal{R}^+$-versal unfolding of $h_{\bm{w}_0}$. 

\item 
$h_{\bm{w}_0}$ has an $A_2$-singularity at $(0,0)$ if and only if $\bm{w}_0$ is the principal normal vector, $f(\bm{0})$ is not an inflection point $($i.e., $a_{2,0} \ne 0)$ and $a_{0,3} \ne 0$. 
When this is the case, $H$ is not an $\mathcal{R}^+$-versal unfolding of $h_{\bm{w}_0}$.

\item 
$h_{\bm{w}_0}$ has an $A_3$-singularity at $(0,0)$ if and only if $\bm{w}_0$ is the principal normal vector, $f(\bm{0})$ is not an inflection point, $a_{0,3} = 0$, and $a_{0,4} a_{2,0} - 3a_{1,2}^2 \ne 0$.
When this is the case, $H$ is not an $\mathcal{R}^+$-versal unfolding of $h_{\bm{w}_0}$.

\item 
$h_{\bm{w}_0}$ has an $A_{\geq 4}$-singularity at $(0,0)$ if and only if $\bm{w}_0$ is the principal normal vector, $f(\bm{0})$ is not an inflection point and $a_{0,3} = a_{0,4} a_{2,0} - 3 a_{1,2}^2 = 0$. 
When this is the case, $H$ is not an $\mathcal{R}^+$-versal unfolding of $h_{\bm{w}_0}$.

\item 
$h_{\bm{w}_0}$ has a $D_4$- or more degenerate singularity at $(0,0)$ if and only if $\bm{w}_0$ is the principal normal vector and $f(\bm{0})$ is an inflection point.
When this is the case, $H$ is not an $\mathcal{R}^+$-versal unfolding of $h_{\bm{w}_0}$.
\end{enumerate}
\end{proposition}

\begin{theorem}
\label{thm:branch}
Let $S$ be a singular surface parameterized by a smooth map-germ $\mathcal{A}$-equivalent to $H_k$, and suppose that the singular point of $S$ is not an inflection point. 
Let $\gamma(t)$ and $\hat\gamma(t)$ be parameterizations, respectively, of the regular and singular branch of the parabolic set on $S$ with $\gamma(0)$ and $\hat\gamma(0)$ being the singular point. 
Let $\bm{b}(t)$ and $\hat{\bm{b}}(t)$ be the unit binormal vectors of $\gamma$ and $\hat\gamma$, respectively. 
Then the height function on $S$ along $\pm\bm{b}(0)$ has an $A_1$-singularity and that along $\pm\lim_{t \to 0^+}\hat{\bm{b}}(t)$ has an $A_2$-singularity at the singular point of $S$.
\end{theorem}

\begin{proof}
We take the parameterization $f$ of $S$ in \eqref{eq:normal_form_2}.
Remark that $a_{2,0}\ne0$ and $a_{0,3}\ne0$. 
Since the parabolic set has a $D_5$-singularity at $(0,0)$, it has a regular branch and a $(2,3)$-cusp branch. 
We may write their parameterizations in the forms
\begin{equation*}
\alpha(t) = \left(t,\, c_1 t + O(t^2)\right),\quad \hat\alpha(t) = (c_2 t^3 + O(t^4),\, \epsilon t^2)\quad(\epsilon=\pm1),
\end{equation*}
where $O(t^n)$ consists of the terms whose degrees are greater than or equal to $n$, respectively. 
Substituting these parameterizations into $P(u,v)=0$ and comparing the lowest-order terms, we obtain 
$$
a_{2,0}(a_{0,3}c_1+a_{1,2})=0,\qquad \dfrac{a_{0,3}}4(4\epsilon a_{2,0} c_2^2 - a_{0,3})=0,
$$
where $\epsilon = 1$ (resp. $-1$) if $a_{0,3}a_{2,0} > 0$ (resp. $< 0$). 
Hence, 
$$
c_1 = -\dfrac{a_{1,2}}{a_{0,3}},\qquad 4\epsilon a_{2,0} c_2^2 - a_{0,3} = 0.
$$
It follows that we can take $\gamma(t)$ and $\hat\gamma(t)$ in the forms
\begin{align*}
\gamma(t) & = f \circ \alpha(t) = \left(t,\, c_1 t^2 + O(t^3),\, \dfrac{a_{2,0}}2t^2 + O(t^3)\right), \\
\hat\gamma(t) & = f \circ \hat\alpha(t) = \left(c_2 t^3 + O(t^4),\, \epsilon c_2 t^5 + O(t^6),\, \dfrac{7\epsilon a_{0,3}}{24} t^6 + O(t^7)\right).
\end{align*}
Straightforward calculations show that
\[
\bm{b}(0) = \left(0,\, -\dfrac{a_{2,0}}{\sqrt{4c_1^2 + a_{2,0}^2}},\, \dfrac{2c_1}{\sqrt{4c_1^2 + a_{2,0}^2}}\right)\ne \pm(0,0,1)
\]
and
\[
\hat{\bm{b}}(t) =\dfrac{\left(\frac{35}4 c_2 a_{0,3}t^8 + O(t^9),\, -\frac{63}4 \epsilon c_2 a_{0,3} t^6 + O(t^7),\, 30 \epsilon c_2^2 t^5 + O(t^6)\right)}{30 c_2^2 |t^5| \sqrt{1 + O(1)}},
 \]
and thus $\lim_{t\to 0^+}\hat{\bm{b}}(t) = (0,0,\epsilon)$.
From (1) and (2) of Proposition \ref{prop:height}, we complete the proof. 
\end{proof}

\end{document}
